\documentclass[10pt]{beamer}
\setbeamertemplate{navigation symbols}{\insertframenumber}

\usepackage[utf8]{inputenc}
\usepackage[T1]{fontenc}
\usepackage[french]{babel}

\usepackage{graphicx}
\usepackage{enotez}

\usepackage{verbatim}
\usepackage{amsmath}
\usepackage{array,multirow,makecell}

\usepackage{listings}
\usepackage{xcolor}

\usepackage{hyperref}

\definecolor{codegreen}{named}{green}
\definecolor{codered}{named}{red}
\definecolor{codeblue}{named}{blue}
\definecolor{codeblack}{named}{black}
\definecolor{backcolour}{rgb}{0.95,0.95,0.92}

\lstdefinestyle{mystyle}{
    backgroundcolor=\color{backcolour},   
    commentstyle=\color{codegreen},
    keywordstyle=\color{codeblue},
    numberstyle=\tiny\color{codeblack},
    stringstyle=\color{codered},
    basicstyle=\ttfamily\footnotesize,
    breakatwhitespace=false,         
    breaklines=true,                 
    captionpos=b,                    
    keepspaces=true,                 
    numbers=left,                    
    numbersep=2pt,                  
    showspaces=false,                
    showstringspaces=false,
    showtabs=false,                  
    tabsize=1
}

\lstset{style=mystyle}
\usepackage[linesnumbered,ruled,french,onelanguage]{algorithm2e}

\makeatletter
\g@addto@macro{\@algocf@init}{\SetKwInput{KwOut}{Sortie}}
\makeatother

\author{Thomas Matthieu, Tellier Basile, Marie Vianney, Siaghi Massinissa}
\date{\today}
\title{Analyse des algorithmes de tri}
\usetheme{Warsaw}

\setlength{\algomargin}{1em}
\begin{document}

\begin{frame}
\maketitle
\end{frame}

\begin{frame}
\frametitle{Table des matières}
\tableofcontents
\end{frame}

\section{Présentation du projet}
\begin{frame}
Objectif du projet :
  \begin{itemize}
      \item <1-> Générateur de données 
      \begin{itemize}
        \item Génération de listes
        \item Gestion du désordre
      \end{itemize}   
      \item <2-> Algorithmes de tri 
      \begin{itemize}
        \item Merge Sort
        \item Counting Sort
        \item Insertion Sort
        \item Bubble Sort 
        \item Quick Sort 
      \end{itemize}
      \item <3-> Analyse des résultats
        \begin{itemize}
          \item Comparaison des algorithmes 
          \item Différences en fonction du pourcentage de désordre 
          \item Différences en fonction du type de désordre 
        \end{itemize}
      \item <4-> Démonstration d'un algorithme de tri
        \begin{itemize}
          \item Démontrer le fonctionnement d'un algorithme 
          \item Comparer les différences de fonctionnement
        \end{itemize}
  \end{itemize}
\end{frame}

\section{Avancement du projet}
\begin{frame}
Avancement actuel du projet :
  \begin{itemize}
    \item <1-> Générateur de données 
      \begin{itemize}
        \item Création d’un tableau trié 
        \item Mélange en fonction de la méthode de désordre sélectionnée 
        \item Renvoi d’un tableau mélangé  
      \end{itemize}
    \item <2-> Algorithmes de tri 
        \begin{itemize}
          \item Implémentation des algorithmes via une classe \texttt{AbstractSort}
          \item Utilisation de sondes pour récupérer le nombre d’accès / swaps / comparaisons
        \end{itemize}
    \item <3-> Analyse des résultats
      \begin{itemize}
        \item Affichage simple des résultats en terminal sans comparaison poussée
      \end{itemize}
    \item <3-> Expérimentation
      \begin{itemize}
        \item Utilisation d’un script shell pour lancer le processus 
        \item Export des données dans un fichier CSV
      \end{itemize}
  \end{itemize}
\end{frame}

\section{Futur du projet}
\begin{frame}
Objectifs pour la suite du projet :
  \begin{itemize}
    \item Interface graphique 
    \begin{itemize}
      \item Visualisation du fonctionnement des algorithmes 
      \item Analyse des performances
    \end{itemize}
    \item Analyse des CSV
    \item Comparaison suivant différents paramètres
    \item Classement général ou avancé
    \item Utilisation d'une graine de génération pour les tableaux
  \end{itemize}
\end{frame}

\begin{frame}
\section{Sources}
  Algorithmes de tri : \url{https://www.geeksforgeeks.org/dsa/sorting-algorithms/}
  \newline
  \url{https://www.youtube.com/watch?v=kPRA0W1kECg}
\end{frame}

\end{document}
